\documentclass[11pt,a4paper]{article}
\usepackage[T1]{fontenc}
\usepackage{lmodern,amsmath,amssymb,amsthm,mathtools,booktabs,microtype}
\usepackage[margin=27mm]{geometry}
\usepackage[colorlinks=true,linkcolor=blue,citecolor=blue,urlcolor=blue]{hyperref}
\newtheorem{theorem}{Theorem}[section]
\newtheorem{lemma}[theorem]{Lemma}
\newtheorem{proposition}[theorem]{Proposition}
\newcommand{\Q}{\mathbb Q}\newcommand{\Z}{\mathbb Z}
\DeclareMathOperator{\Tr}{Tr}\DeclareMathOperator{\Nm}{N}
\DeclareMathOperator{\rad}{rad}
\setlength{\emergencystretch}{2em}
\title{A Quartic in Ninety-Eight Variables under the Rational $abc$ Conjecture}
\author{}\date{9 November 2026}
\begin{document}\maketitle
\begin{abstract}
Assuming the rational $abc$ conjecture and the specified elliptic-group,
pole-realization and simultaneous prime-pattern inputs, we construct a
family of integer quartics in ninety-eight total variables. Every
sufficiently large fixed odd input shift gives rational input projection
exactly $\Z$. The quartic is a sum of eighty-six primitive quadratic
squares and has exactly $882$ monomials. Its coefficient height, for a
shift $B\ge3$, is $\max(11600064,2B^4+169997)$. Six rational contacts
symmetric about one permit a short quadratic evaluation and normalization.
A height estimate for a sum of sixth powers bounds the input together
with one elliptic index. All six local tests have order one; only plain
elliptic formal parameters are used. The sufficient shift is ineffective.
\end{abstract}

\section{Statement and arithmetic assumptions}

For a reduced fraction $u/v$, with $v>0$, put
$h(u/v)=\log\max(|u|,v)$. Write $h(1:u:v)$ for the logarithmic
projective height of the indicated rational point. The coefficient height
of an integer polynomial is the maximum absolute value of its coefficients.
All variable counts include the distinguished input.

\begin{theorem}\label{thm:main}
Under the assumptions below, there is a computable family
\[
 P_B\in\Z[X,Y_1,\ldots,Y_{97}],\qquad B\ge3\text{ odd},
\]
such that every member has degree exactly four, content one, exactly
$882$ monomials, and coefficient height
\[
 H(P_B)=\max(11600064,2B^4+169997).
\]
There is a constant $B_*$ such that, for every odd $B\ge B_*$,
\[
 x\in\Z\quad\Longleftrightarrow\quad
 \exists\boldsymbol y\in\Q^{97},\quad P_B(x,\boldsymbol y)=0.
\]
The integer $B$ is a fixed coefficient parameter. No numerical sufficient
value of $B_*$ is determined. In particular, the coefficient fixture
$B=3$ is not asserted to define the integers.
\end{theorem}

Let $K=\Q(\rho)$, where $\rho^2=2$, let $\sigma$ be conjugation, and put
\[
 d=3-\rho,\quad D=d^2=11-6\rho,\qquad E:y^2=x^3-Dx,
\]
\[
 R_0=(1+\rho,2\rho),\qquad
 P=2R_0=(x_P,y_P)
   =\left(\frac{57-28\rho}{8},\frac{-660+417\rho}{32}\right).
\]
The arithmetic assumptions have the following precise scope.
\begin{enumerate}
\item The rational $abc$ conjecture: for every $\epsilon>0$ there is
$K_\epsilon\ge1$ such that $c\le K_\epsilon\rad(abc)^{1+\epsilon}$
for pairwise coprime positive integers $a+b=c$.
\item The elliptic source gives $2E(K)=\Z P$, with $P$ of infinite
order, and the index-height comparison stated in~\cite{Dependencies}.
The square-abscissa characterization proved in~\cite{Squares} then
applies to all nonzero multiples of $P$.
\item For a nonzero rational $b$, set $B_0=17b$. The pole equations
\begin{equation}\label{eq:originalpole}
 \begin{split}
 p_0^2&=B_0(1-e_0)(1+3e_0),\\
 B_0q_0^2&=(s_0^2-B_0e_0(B_0-b^2))
                  (s_0^2+3B_0e_0(B_0-b^2))
 \end{split}
\end{equation}
exclude negative odd valuations of $b$ outside $\{2,3,17\}$.
For inputs satisfying
\begin{equation}\label{eq:prepared}
 \begin{gathered}
 0<b<17,\qquad b\text{ has square reduced denominator},\\
 b\in\Q_2^{\times2}\cap\Q_3^{\times2}\cap\Q_{17}^{\times2},
 \qquad v_2(b)\ge3,\quad v_3(b)\ge1,\quad v_{17}(b)\ge2,
 \end{gathered}
\end{equation}
the retained restricted-converse and descent hypotheses give a witness
with $p_0q_0s_0\ne0$. The hypotheses and their nonzero-witness consequence
are specified in~\cite{Dependencies,Source27,PoleDensity}.
\item The simultaneous prime-pattern input has exactly the quantified
form in~\cite{Dependencies,Selection}: arbitrary fixed rational contact
coefficients after denominator clearing, the prescribed Gaussian
congruences and conjugate reduction orders, fixed avoidance and
divisibility, and compatible real approximation. Its analytic source is
Sun's specified Theorem~4.1 and the estimates pinned in those papers.
All scales and finite preparation data are fixed before prime selection.
\end{enumerate}
The exponent-nineteen contact-height input is not used here.

The one-point case of the proved aggregate trace-height theorem
\cite{Trace}, obtained from von Buhren's family height theorem, gives
constants $A>0$, $J\ge0$ such that
\begin{equation}\label{eq:traceheight}
 n^2\le A h(T(n))+J,\qquad
 T(n)=\Tr x([n]P),\quad n\in\Z\setminus\{0\}.
\end{equation}
The additive constant is ineffective. Each abscissa is a positive square
at both real embeddings and exceeds $d$ there, so $T(n)>6$.

\section{Two height consequences of $abc$}

Set
\[
 \Theta_6=\{1\pm1,1\pm2,1\pm3\}
          =\{-2,-1,0,2,3,4\}.
\]
For $t\notin\Theta_6$, let $M(t)$ be the product of primes $p>76176$
for which $v_p(t-\theta)>0$ is odd for some $\theta\in\Theta_6$.

\begin{lemma}\label{lem:contact}
Assume rational $abc$. There is a fixed $C_0\ge0$ such that
\[
                         h(t)\le\frac23\log M(t)+C_0.
\]
\end{lemma}
\begin{proof}
For a primitive pair $(u,v)$, $v>0$, use the squarefree binary form
$G(u,v)=\prod_{\theta\in\Theta_6}(u-\theta v)$ of degree six.
Granville's consequence of $abc$~\cite[Proposition~1]{Granville},
with $\epsilon=1/4$, gives
\[
 \log\rad(G(u,v))\ge\frac{15}{4}\log\max(|u|,v)-C.
\]
For a nonzero integer $m$, let $\operatorname{sf}(m)$ be the product
of primes having odd valuation in $m$. Prime by prime,
$\rad(m)^2\le |m|\operatorname{sf}(m)$.
Also $|G(u,v)|\le C'\max(|u|,v)^6$.
Outside the fixed cutoff, no prime divides two distinct factors of
$G$, and a prime dividing $v$ divides none of them. Its squarefree
part outside the cutoff is therefore precisely $M(u/v)$.
The omitted squarefree factors are bounded by a fixed product of primes.
Consequently
\[
 \log M(u/v)\ge\left(2\cdot\frac{15}{4}-6\right)h(u/v)-C''
                 =\frac32h(u/v)-C'',
\]
as required.
\end{proof}

\begin{lemma}\label{lem:pair}
Assume rational $abc$. There is a fixed $C_1\ge0$ such that
\[
                    h(1:u:v)\le\frac25h(u^6+v^6)+C_1
                    \qquad(u,v\in\Q).
\]
One may take $C_1=\frac7{30}\log2+\frac{11}{30}\log K_{1/11}$.
\end{lemma}
\begin{proof}
The zero-coordinate cases follow from $h(w^6)=6h(w)$.
Otherwise write $u=gA/D_0$, $v=gB/D_0$, where
$\gcd(A,B)=\gcd(g,D_0)=1$, $g,D_0>0$. Put
$H=\max(|A|,|B|)$ and $N=A^6+B^6$, so $H^6\le N\le2H^6$.
Let $k=\gcd(N,D_0^6)$, $m_0=N/k$, $n_0=D_0^6/k$.
Then $u^6+v^6=g^6m_0/n_0$ is reduced. Write
$M=\max(g^6m_0,n_0)$.

Every prime dividing $N$ either divides $D_0$ or survives in $m_0$,
so $\rad(N)\le D_0m_0$. Apply $abc$ with $\epsilon=1/11$ to
$A^6+B^6=N$. The three positive terms are pairwise coprime, and
\[
 H^{7/2}\le K_{1/11}^{11/12}D_0m_0.
\]
Using $D_0^6=n_0N/m_0\le2n_0H^6/m_0$ gives
\[
 H^{5/2}\le2^{1/6}K_{1/11}^{11/12}m_0^{5/6}n_0^{1/6}.
\]
Moreover $\max(gH,D_0)\le2^{1/6}H(M/m_0)^{1/6}$.
Taking logarithms and combining yields
\[
 \begin{split}
 h(1:u:v)&=\log\max(gH,D_0)\\
 &\le C_1+\frac16\log m_0+\frac1{15}\log n_0+\frac16\log M
 \le C_1+\frac25\log M.
 \end{split}
\]
This includes common numerator factors and all denominator cancellation.
\end{proof}

\section{Quadratic point equations}

Put $\beta=2X^2+B$. If $X=a/b$ is reduced, integrality of $\beta$
implies $b^2\mid2$, hence $b=1$. Thus $\beta\in\Z$ if and only if
$X\in\Z$; for integer input $\beta$ is odd, and always $\beta\ge B$.

All equations in this section are over $K$ and mean their two rational
component equations. Products of stored quantities have degree at most two.
For the root point retain $h,j,a\in K$ and impose
\begin{equation}\label{eq:root}
                       a=h^2,\qquad j^2=a^2-D.
\end{equation}
Then $A=(a,hj)=[n]P$ for some nonzero integer $n$ by the
square-abscissa theorem~\cite{Squares}. Retain $\chi,f,S\in K$ with
\begin{equation}\label{eq:H}
                    \chi f=1,\qquad S^2=\chi-Df.
\end{equation}
They give $H=(\chi,\chi S)\in E(K)$. Write
$H=[m]R_0+T_0$, $T_0\in E[2](K)$, and put $h_*=m/2$.
This is an index notation; it does not mean multiplication of a point
by a noninteger rational number.

\subsection{A pair of points with fixed difference}

Retain $V,Q,W_A,I_A,K_A\in K$ satisfying
\begin{equation}\label{eq:pairchart}
 \begin{split}
 hI_A&=(j+DW_A)V+aQ,\\
 hK_A&=-(j+DW_A)Q-aV,\\
 VQ&=aW_A,\\
 V^2+Q^2+2jW_A+DW_A^2&=1.
 \end{split}
\end{equation}

\begin{lemma}\label{lem:pairchart}
Equations~\eqref{eq:root} and~\eqref{eq:pairchart} describe points
$U,V_1\in\Z P$ with $U-V_1=A$, allowing either numerator point
to be the identity. Both $I_A$ and $K_A$ are nonzero. On the
nonidentity branches,
\[
 U=\left(\frac a{V^2},\frac{aI_A}{V^3}\right),\qquad
 V_1=\left(\frac a{Q^2},\frac{aK_A}{Q^3}\right),
 \qquad z(U)=-\frac V{I_A},\quad z(V_1)=-\frac Q{K_A},
\]
where $z=-x/y$. In all branches the quantities $R_U,R_1\in K$
defined by
\begin{equation}\label{eq:recipsquares}
             I_A^2=a-DaR_U,\qquad K_A^2=a-DaR_1
\end{equation}
are $f(U)^2,f(V_1)^2$, with $f=1/x$ extended by zero at the identity.
\end{lemma}
\begin{proof}
The square chart gives $haj\ne0$.
Squaring the first two equations and using the last two gives
$aI_A^2=a^2-DV^4$ and $aK_A^2=a^2-DQ^4$.
For $VQ\ne0$ these are exactly the two curve equations. Their
abscissas are squares, so the points lie in $\Z P$.
The line with slope $(1/W_A-j)/h$ through $U$ also passes through
$-V_1$ and $-A$. Substitution gives its three abscissas
$a/V^2,a/Q^2,a$, counted with multiplicity. The chord-and-tangent
law therefore gives $U-V_1=A$, including the tangent case.
Nontorsion of the nonidentity points implies $I_AK_A\ne0$.

If $V=0$, then $W_A=0$, $Q^2=1$, $I_A=hQ$ and
$K_A=-jQ/h$. The interpretation is $U=O,V_1=-A$.
If $Q=0$, then $W_A=0$, $V^2=1$, $I_A=jV/h$ and
$K_A=-hV$, giving $U=A,V_1=O$. The two zero cases cannot occur
together. The asserted nonvanishing and reciprocal-square identities
also hold there.

Conversely, let $U,V_1\in\Z P$ be nonidentity points with $U-V_1=A$.
Their abscissas are squares. The line through $U,-V_1,-A$ has equation
$y=\mu x+c$, with $c=-h(\mu h+j)$. Put $R=\mu h+j$.
Vieta's identities give
$x(U)x(V_1)=R^2$ and $x(U)+x(V_1)=\mu^2-a$; in particular $R\ne0$.
Choose square roots $u_0^2=x(U)$, $v_0^2=x(V_1)$ with $u_0v_0=R$,
and put
\[
 V=h/u_0,\qquad Q=h/v_0,\qquad W_A=1/R,\qquad
 I_A=y(U)V^3/a,\qquad K_A=y(V_1)Q^3/a.
\]
The two Vieta identities give the last two chart equations, and the
line ordinates give the first two. This also covers tangent lines.
The identity branches are supplied by the explicit formulas above.
\end{proof}

\subsection{Subtraction of the fixed point}

Retain $g,F_0,z_1\in K$ and impose
\begin{equation}\label{eq:fixedsubtract}
 \begin{split}
 8(x_P\chi+D)g+(39-36\rho)+(8-14\rho)S&=0,\\
 fF_0&=g^2,\qquad K_Az_1+Q=0.
 \end{split}
\end{equation}
On the admissible domain, $F_0=f(H-P)$ and $z_1=z(V_1)$.
Define the affine expressions
\begin{equation}\label{eq:ND}
 \begin{split}
 N&=(-312-69\rho)+(-236-120\rho)S,\\
 D_0&=(-576-256\rho)+(312+69\rho)S+(504-1344\rho)g.
 \end{split}
\end{equation}
Retain $L,L_0\in K$ with
\begin{equation}\label{eq:ratios}
                         I_AL=VS,\qquad L_0N=z_1D_0.
\end{equation}
After the exclusions proved below, these mean
\[
 L=\frac{z(U)}{z(H)},\qquad
 L'=\sigma L_0=\sigma\!\left(\frac{z(V_1)}{z(H-P)}\right).
\]
No coordinate for $z(H-P)$ is retained.

For clarity, the fixed subtraction formulas follow from the usual
addition slope $(\chi S+y_P)/(\chi-x_P)$.
Direct reduction using $S^2=\chi-D/\chi$ gives
\[
 f(H-P)=\chi g^2,\qquad z(H-P)=N/D_0.
\]
Before multiplication by the common unit $3+2\rho$, the numerator
is $32(y_P-x_PS)$ and the denominator is
$(-704+384\rho)+(660-417\rho)S+(6888-5040\rho)g$.
The numerator vanishes only when
$S=y_P/x_P$, and the curve equation then gives
\[
                    (\chi-x_P)(\chi+D/x_P)=0.
\]
The latter alternative includes $H=-P+(0,0)$ and must be excluded,
along with $H=P$, before division. An exact function-field check of
these identities is included with the coefficient certificate.

\section{Six tests and two comparison scales}

Set
\[
 C=L-L',\qquad I=L-\beta,\qquad T=\Tr a,
 \qquad a_0=1,\quad a_1=\beta^6+T^6.
\]
The six field tests are $C,I,f(H),f(U),\sigma f(H-P),\sigma f(V_1)$.
Retain the four reciprocal squares, $C^2,I^2$, and the fourth powers
$f(U)^4,\sigma f(V_1)^4$ as quadratic outputs. The two fields in
\eqref{eq:recipsquares} supply two reciprocal squares directly.
Four scalar outputs compute the comparison scale without taking any
root or making any division:
\[
 s_h=\beta^2+T^2,\quad p_h=\beta T,\quad
 v_h=s_h^2-3p_h^2,\quad a_1=s_hv_h=\beta^6+T^6.
\]

Retain a rational $\lambda$. For $a\in\{a_0,a_1\}$ put $t=a\lambda$,
$\delta=t-1$, and use scalar coordinates $r,q,c,k$ satisfying
\begin{equation}\label{eq:contactcircuit}
 r=\delta^2,\qquad c=qr,\qquad k=cr,\qquad
 kr-14k+49c-36q=r^2.
\end{equation}
Define, without storing another coordinate,
\[
 F=(r-1)(r-4)(r-9)=r^3-14r^2+49r-36
   =\prod_{\theta\in\Theta_6}(t-\theta).
\]
The displayed quadratic equations imply exactly $Fq=r^2$ and $k=qr^2$.
They compute the needed normalization without a coordinate for $F$.
Store $q^2,q^4$ and the two field factors
\begin{equation}\label{eq:packed}
 \begin{split}
 T_1&=1+qC^2+q^2f(H)^2+q^4f(U)^4,\\
 T_2&=1+qI^2+q^2\sigma f(H-P)^2+q^4\sigma f(V_1)^4,
 \qquad W=T_1T_2=w_0+w_1\rho.
 \end{split}
\end{equation}
Finally retain scalars $\tau,b$ with
\begin{equation}\label{eq:trace}
             w_0^2-2w_0\tau-2w_1^2=0,\qquad
             kb=\tau^2.
\end{equation}

The normalized pole uses seven further rational coordinates
$e,u,p,v,w,p_2,w_2$ and five quadratic equations
\begin{equation}\label{eq:normalizedpole}
 \begin{gathered}
 p_2=p^2,\qquad 17bp_2=289+34e-3e^2,\\
 e(17-b)=17uv,\qquad w_2=w^2,\qquad
 17bw_2=u^2+34uv-867v^2.
 \end{gathered}
\end{equation}
For $b\ne0$ these give~\eqref{eq:originalpole} under
\[
 e_0=e/17,\quad p_0=bp,\quad s_0=bu,\quad q_0=b^2uw.
\]
Conversely a witness with $s_0\ne0$ gives the displayed normalized
coordinates by these formulas and $v=e(17-b)/(17u)$.
This explains both the pole obstruction and its use in completeness.

\begin{lemma}\label{lem:domain}
Every rational solution of the quadratic source has $\delta Fqk\ne0$
in both detectors. Consequently
\[
 q=\frac{\delta^4}{F},\qquad
 k=\frac{\delta^8}F=\frac1R,
 \qquad R=\frac F{\delta^8}.
\]
\end{lemma}
\begin{proof}
If $F=0$, then $r^2=\delta^4=0$, whereas at $t=1$ the contact
product equals $-36\ne0$.
If $\delta=0$, the equations give $q=0$, $T_1=T_2=W=1$,
$\tau=1/2$, and $k=0$, contradicting $kb=\tau^2$.
The remaining assertions follow directly. In particular $\lambda\ne0$
because zero is a contact.
\end{proof}

\section{Local extraction and denominator exclusions}

At a prime $p>76176$, the curve has good reduction and $D$ is a unit.
For a nonzero finite reciprocal abscissa $f(Q)$, its valuation at
either place of $K$ above $p$ is even. If $v(x)<0$, the curve equation
gives $2v(y)=3v(x)$; if $v(x)>0$, it gives $2v(y)=v(x)$.
The unit case is immediate. Zero reciprocal abscissas are interpreted
as having infinite valuation.
Before the exceptional $H$ classes are excluded, the coordinate $F_0$
need not yet be interpreted as $f(H-P)$. Its even valuation nevertheless
follows directly from $fF_0=g^2$ and the even valuation of $f=f(H)$;
$F_0=0$ is allowed. Thus the valuation argument below applies to the
literal equations before any division in the fixed subtraction chart.

Let $p\mid M(t)$, put $\nu=v_p(R)$, and normalize valuations by
$v(p)=1$. Then $\nu$ is positive and odd, $\delta$ is a unit, and
$v(q)=-\nu$. For each factor in~\eqref{eq:packed}, the three
nonconstant finite valuations belong respectively to
\[
                     1+2\Z,\qquad 2+4\Z,\qquad4+8\Z.
\]
Thus all term valuations, including zero for the constant term, are
distinct. Each factor is nonzero and has valuation at most zero;
valuation zero implies residue one. A zero test merely removes a term.

\begin{proposition}\label{prop:local}
At every contact prime, there is one place at which all six tests have
positive valuation. The four reciprocal-abscissa tests have valuation
at least two.
\end{proposition}
\begin{proof}
The preceding argument gives $W\ne0$. Equation~\eqref{eq:trace} then
implies $w_0\ne0$, since $w_0=0$ would imply $w_1=0$. Hence
\[
 \tau=\frac{\Nm(W)}{2w_0}=\frac1{\Tr(W^{-1})}\ne0,
 \qquad b=R\tau^2\ne0.
\]
Put $\ell=\max_{v\mid p}v(W)\le0$. The trace inequality gives
$v_p(b)\le\nu+2\ell$. This valuation is odd, so the pole obstruction
forces $0<v_p(b)\le\nu+2\ell$, hence $\ell>-\nu/2$.
At a place attaining $\ell$, each factor valuation is at least $\ell$.
Each of its terms therefore has valuation greater than $-\nu/2$.
For the three columns this gives respectively
\[
 v(C),v(I)>\nu/4,\qquad
 v(f(H)),v(\sigma f(H-P))>3\nu/4,\qquad
 v(f(U)),v(\sigma f(V_1))>7\nu/8.
\]
Integer rounding and evenness give the assertion.
\end{proof}

The finite $H$ chart may contain torsion translates. Before interpreting
the subtraction formulas or either ratio, exclude
\[
                  E[2](K)\ \cup\ (P+E[2](K))\ \cup\ (-P+E[2](K)).
\]
Every permitted point in this finite set has reciprocal abscissa a unit
above $p>76176$. The torsion translation formulas, together with
\[
 \Nm(x_P)=41^2/64,\quad
 \Nm(x_P-d)=17^2/64,\quad \Nm(x_P+d)=63^2/64,
\]
give this directly. If $H$ belonged to the set, Proposition~\ref{prop:local}
would force both contact radicals to equal one. Lemmas~\ref{lem:contact}
and~\ref{lem:pair}, and $h(a_1)\le h(a_1\lambda)+h(\lambda)$, would
then give $h(1:\beta:T)\le\frac45C_0+C_1$. Since $\beta\ge B$, this is
impossible for every sufficiently large fixed $B$.

Consequently $h_*\notin\{-1,0,1\}$ and all required subtraction
denominators are nonzero. In particular $x_P\chi+D\ne0$, $N\ne0$,
$D_0\ne0$, and $z(H),z(H-P)\ne0$. The formulas above have their
stated point interpretation. This exclusion does not use a logarithmic
ratio and therefore precedes, rather than assumes, its validity.

\section{Soundness for all sufficiently large shifts}

Write $U=[u]P$, so $V_1=[u-n]P$. At a contact place supplied by
Proposition~\ref{prop:local}, $H,U$ lie in the primary formal group
and $H-P,V_1$ in the conjugate formal group. The formal parameter satisfies
\begin{equation}\label{eq:log}
 \log_E(Q)=z(Q)\mathcal A(Q),\qquad
 \mathcal A(Q)\equiv1\pmod{p^{4v(z(Q))}}.
\end{equation}
This is the constant-term case of the proved formal-logarithm estimate
in~\cite{Logs}. The unit is defined to be one at the identity.

The integral group identities
\[
                [m]U=[2u]H,\qquad
                [m-2]V_1=[2(u-n)](H-P)
\]
kill the torsion component of $H$. Their full logarithmic ratios are
$r=u/h_*$ and $r'=(u-n)/(h_*-1)$. At least one of $h_*,h_*-1$ is a
$p$-adic unit, so at least one full ratio is integral. The units in
\eqref{eq:log} preserve the corresponding approximate valuation, and
the positive-order test $C=L-L'$ transfers integrality to the other
approximate ratio and hence to both full ratios.
Only after this step does~\eqref{eq:log} give absolute congruences
$L\equiv r$ and $L'\equiv r'\pmod p$. Therefore
\[
 n=h_*r-(h_*-1)r'\equiv L\equiv\beta\pmod p.
\]
The argument also covers identity numerator points and half-integral
$h_*$, since $p$ is odd. In particular $\beta$ is integral at $p$.

If $\beta\ne n$, let $N_\beta$ be the positive absolute reduced
numerator of $\beta-n$. Both $M(\lambda)$ and $M(a_1\lambda)$ divide
$N_\beta$. Lemma~\ref{lem:contact} gives
\[
 h(a_1)\le h(\lambda)+h(a_1\lambda)
               \le\frac43\log N_\beta+2C_0.
\]
Put $H_2=h(1:\beta:T)$. Lemma~\ref{lem:pair} implies
\[
 H_2\le\frac8{15}\log N_\beta+C_2,\qquad C_2=\frac45C_0+C_1.
\]
For a reduced fraction $\beta=a/b$,
$|a-nb|\le(1+|n|)\max(|a|,b)$. Since $h(\beta)\le H_2$,
\[
 \log N_\beta\le H_2+\log(1+|n|),\qquad
 H_2\le\frac87\log(1+|n|)+\frac{15}{7}C_2.
\]
The trace-height bound~\eqref{eq:traceheight} now yields
\[
 n^2\le\frac{8A}{7}\log(1+|n|)+\frac{15A}{7}C_2+J.
\]
Its right side grows logarithmically, so $|n|$ and then $H_2$ are
uniformly bounded. Every constant here is independent of $B$.
Since $\beta\ge B$, choosing one sufficiently large odd shift excludes
$\beta\ne n$. Thus $\beta=n\in\Z$, and consequently $X\in\Z$.
Together with the earlier exceptional-point exclusion this proves
soundness for every odd $B$ beyond one fixed threshold.

\section{Simultaneous completeness}

Fix $X\in\Z$ and an odd $B\ge3$. Put $n=\beta$ and $A=[n]P$;
this fixes $T$ and both positive rational scales. Choose a nonzero
reference integer divisible by $64$ and compact real and finite-place
neighborhoods for the points
\[
 H=[h]P,\quad U=[hn]P,\quad H-P=[h-1]P,\quad
 V_1=[(h-1)n]P
\]
which avoid the finitely many chart denominators. In these neighborhoods
every rational test has a fixed lower valuation bound at the finitely
many preparation primes. All these choices precede prime selection.

Choose a fixed $D_1>0$ clearing $D_1^2\theta/a$ for both scales and
all contacts, and sufficiently divisible at the preparation primes.
Apply the retained prime-pattern input to the distinct coefficients
$-D_1^2\theta/a$, with repetitions removed, and write
\[
                \lambda=\frac{u}{D_1^2w^2},\qquad L_1=D_1^2w^2.
\]
Its prime values are $u-(\theta/a)L_1$. The zero contact makes $u$
one of the selected primes. The supplied conjugate reduction-order
conditions and the compatible congruence progression in~\cite{Dependencies}
give $h$ for which $H,U$ are formal at one place above each selected
prime and $H-P,V_1$ are formal at the conjugate place. The fixed
congruence and real neighborhoods are retained. Only natural depth one
is requested at either branch.

Both full logarithmic ratios equal $n$. Hence the four reciprocal
tests have valuation at least two and $C,I$ have valuation at least
four at the chosen place. Since $v_p(q)=-1$, every nonconstant term
of both factors in~\eqref{eq:packed} has positive valuation there.
At the other place each factor is either a unit of residue one or has
negative valuation. Thus $\Tr(W^{-1})$ has residue one or two and is
a unit. The pole input has valuation one at every selected contact prime.

For arbitrary positive rational $a$ the exact identity
\begin{equation}\label{eq:primeproduct}
 R=\frac{a^6\prod_{\theta\in\Theta_6}(u-(\theta/a)L_1)}
             {L_1^6(a\lambda-1)^8}
\end{equation}
has only square factors apart from the selected prime numerators.
In particular no square condition on $a$ is needed. At every other
prime $v_p(R)$ is even, and the reciprocal trace changes valuations
by an even number. These facts give a square reduced denominator for $b$.

At a preparation prime, increasing the fixed denominator divisibility
makes $v_p(t)\ll0$. Then
\[
 v_p(q)=v_p(R)=-2v_p(t),\qquad
 R=\delta^{-2}\prod_{\theta\in\Theta_6}
                   \left(1+\frac{1-\theta}{\delta}\right).
\]
The product is arbitrarily close to one and is a local square at
$2,3,17$. All weighted test terms are small, so $\Tr(W^{-1})$ is
arbitrarily close to two. Sufficient divisibility therefore gives the
three local squareclasses and valuation depths in~\eqref{eq:prepared}.

On the fixed compact real neighborhoods, as the lower bound on $\lambda$
tends to infinity, $R\sim q\sim t^{-2}$, $W\to1$, and
$b\sim1/(4t^2)$. Both real embeddings of every factor are positive
for sufficiently large $\lambda$. Hence $0<b<17$ and every required
denominator is nonzero. These bounds are uniform on the chosen compact
sets, so they can be fixed before applying the prime-pattern theorem.
Each prepared $b$ has a nonzero pole witness, which converts to
\eqref{eq:normalizedpole}. All quadratic coordinates can now be filled
by their stated rational values. This proves completeness for every
integer input and every odd $B\ge3$.

\section{The polynomial and its exact certificate}

Each field value represents two rational coordinates. For the abscissas
$a,\chi$, the four reciprocal squares and $C^2$, the stored second
coordinate is one half of the $\rho$ coefficient. Other stored field
values use their ordinary coefficients. These rational changes of
coordinates are invertible; they are fixed in the coefficient tables.
Each component equation is cleared to a primitive integer quadratic,
with positive first nonzero coefficient in the prescribed monomial order.
Their sum of squares defines $P_B$.

The coordinate count is
\begin{center}
\begin{tabular}{lr}
\toprule
Coordinates & Number\\\midrule
Input, its square, and the common slope &3\\
Root chart for $A$ &6\\
Chart for $H$ &6\\
Fixed subtraction quantities $g,F_0$ &4\\
Pair chart for $U,V_1$ &10\\
$z(V_1)$ and the two ratios &6\\
Four reciprocal squares and two comparison squares &12\\
Two fourth powers &4\\
Four scalar comparison-scale outputs &4\\
Two detector blocks, including their pole coordinates &42\\
The second argument $a_1\lambda$ &1\\\midrule
Total &98\\\bottomrule
\end{tabular}
\end{center}
Each detector uses two contact outputs $r,c$, three weight coordinates,
four factor components, two product components, the three scalars
$\tau,k,b$, and seven pole coordinates: twenty-one in all.
There are $86$ primitive quadratic equations, so over $\Q$ their
sum of squares vanishes exactly when they all vanish.

The complete specification is in
\nolinkurl{computation/quartic/family_abc98}. The builder and independent
checker reconstruct all $86$ rows and every coefficient of the full
quartic. There are $885$ terms when powers of $B$ are stored separately,
and $882$ ordinary monomials for every positive integer $B$.
The coefficient of $X^4$ is one, proving degree four and content one.
Only three ordinary coefficients contain more than one power of $B$:
the constant coefficient $2B^4+169997$, the coefficient
$48B^2+1$ of the square of the input-square coordinate, and $12B^2-19$
at the square of the rational component of the root abscissa.
All remaining nonconstant coefficient profiles are single powers of $B$;
their maximal multipliers at powers one, two and three are respectively
$64,48,16$. The largest parameter-free coefficient has absolute value
$11600064$. It follows directly that
\[
 H(P_B)=\max(11600064,2B^4+169997)
       =2B^4+169997\quad\text{for odd }B\ge49.
\]
No coefficient cancellation changes the stated support for positive integer $B$.

The independent checker imports no producer. It also reconstructs
the fixed-subtraction compression and both normalizations. A separate
PARI computation verifies the fixed subtraction identities in the exact
function field $K(\chi)[S]/(S^2-\chi+D/\chi)$ against the elliptic
addition law. These are finite algebraic certificates. The arithmetic
assumptions and the ineffective sufficient shift remain those stated
in Section~1.

\begin{thebibliography}{10}
\bibitem{Granville}
A.~Granville, Rational and integral points on quadratic twists of a given
hyperelliptic curve, \emph{International Mathematics Research Notices}
(2007), article rnm027, Proposition~1.
\url{https://doi.org/10.1093/imrn/rnm027}; author's manuscript,
\url{https://dms.umontreal.ca/~andrew/PDF/hyperelliptics.pdf}.
\bibitem{Dependencies}
\emph{Arithmetic dependencies of the construction with 29 variables},
companion paper, \nolinkurl{papers/arithmetic_dependencies}.
\bibitem{Selection}
\emph{Separation and denominator conditions in quadratic prime patterns},
companion paper, \nolinkurl{papers/prime_selection}.
\bibitem{Squares}
\emph{Square abscissas and direct charts for the elliptic subgroup},
companion paper, \nolinkurl{papers/square_abscissa}.
\bibitem{Trace}
\emph{Aggregate trace heights of elliptic multiples}, companion paper,
\nolinkurl{papers/trace_heights}; it applies J.~von Buhren,
\emph{J. Th\'eor. Nombres Bordeaux} \textbf{29} (2017), 307--320,
Theorem~1.1, \url{https://www.numdam.org/item/10.5802/jtnb.981.pdf}.
\bibitem{Logs}
\emph{Uniform formal-logarithm estimates}, companion paper,
\nolinkurl{papers/elliptic_logarithms}.
\bibitem{Source27}
\emph{A defining polynomial in twenty-seven variables: an existence
theorem with a uniform degree bound}, companion paper,
\nolinkurl{papers/few_variables27}, nonzero pole realization lemma.
\bibitem{PoleDensity}
\emph{Nonzero witnesses in the pole equations}, companion paper,
\nolinkurl{papers/pole_parameters/pole_witness_density_20261109_v14.tex}.
\end{thebibliography}
\end{document}
